\section{What the reviewed studies report}
\label{sec:findings}


\begin{table*}[!tb]
\centering\footnotesize\setlength\tabcolsep{3pt}\renewcommand{\arraystretch}{1.08}
\caption{What this corpus reports, and what it does not: the five conclusions, the statements of
\S\ref{sec:findings} behind each (under the conclusion), the claim each stops short of, the property that limits it most, and the
experiment of \S\ref{sec:future} that would test it. Counts are on all records.}
\label{tab:synthesis}
\begin{tabular}{@{}L{2.45cm}L{3.3cm}L{3.2cm}L{3.2cm}L{3.2cm}@{}}
\toprule
Conclusion (rows) & Reported in this corpus & Not established & Principal limitation & Experiment that would test it \\
\midrule
\textbf{Availability against use} \newline {\scriptsize F1a, F1a$'$, F1b, F1c, F2, F2b, F9a} & findings, anatomy, demographic and acquisition attributes recoverable from frozen encoders in four to seven modalities & that the output uses them: matched control in \Vspca{}--\Vspc{} of \Vspcn{} targeted interventions; no concept-specific causal contribution by the criterion of \S\ref{sec:causal-claims} & breadth is breadth over records; the safeguards pooled into one control are not equivalent & probe paired with a matched intervention or erasure at the same locus and end-point (item 2) \\
\addlinespace[2pt]
\textbf{Acquisition and site structure} \newline {\scriptsize F2b, F3a--F3c} & centre, scanner and acquisition structure recoverable from pathology and brain-MRI embeddings & that removing it restores generalisation: the downstream effect is operation-specific & F3b rests on five records from five studies, two of them flagged by the defect screen; pathology dominates & downstream end-point, not only invariance, after harmonisation on a held-out site \\
\addlinespace[2pt]
\textbf{Meaning of discovered features} \newline {\scriptsize F7} & \Nsaeval{} of \Nsae{} dictionary studies report a semantic check of feature meaning against clinicians or labels, all encoder-side & cross-seed identity of features or names; completeness or causal meaning; any decoder-residual dictionary & the row selects studies that report a check: how it was done, not how often it succeeds & two differently seeded dictionaries on one encoder; matched-feature fraction; names validated on region annotations (item 3) \\
\addlinespace[2pt]
\textbf{Medical against general pre-training} \newline {\scriptsize F10, F11} & in same-data comparisons, selected geometric and shortcut properties reported in both kinds of model & an effect of medical pre-training: better decodability appears only in uncontrolled comparisons & \Nmvgrec{} comparisons from \Nmvgstud{} studies vary architecture, objective, scale and data together; most strata hold one to five records & paired pre-training on matched corpora with objective, architecture and compute fixed, or a medical-fraction mixture read as a dose--response \\
\addlinespace[2pt]
\textbf{Where to measure} \newline {\scriptsize F8b, \S\ref{sec:loci}} & connector output and text encoder almost unmeasured as loci of their own & that decodability depends on readout and locus; any ordering of model, depth and pooling & readout evidence: three studies varying different things; locus evidence: one uncontrolled successive-locus probe & one finding probed across encoder, connector output and every language-model layer of one report generator (item 1) \\
\bottomrule
\end{tabular}
\end{table*}

% the evidence map of \S\ref{sec:findings}: an upright two-column float
\begin{figure*}[!tbp]
\centering
\includegraphics[width=\linewidth]{fig6_findings.pdf}
\caption{The evidence map: for each statement of \S\ref{sec:findings}, the controls its records report, in
five groups, as $k/n$ over the records whose design can carry the control. \emph{Chance calibration}: a
label-permutation null or a reference distribution. \emph{Attribution to training}: a randomly initialised
encoder, a planted ground truth or a matched object. \emph{Comparison model}: the same statistic on
another model. \emph{Perturbation test}: an occlusion, patching or counterfactual check of a localisation.
\emph{Matched control}: a control direction or unit for a targeted intervention. A dash marks a control
the row's designs cannot carry; F7 has no cell because it selects on semantic validation. The
\Vnclasses{} underlying classes, and what each licenses, are listed in the supplement.}
\label{fig:findings}
\end{figure*}




Table~\ref{tab:synthesis} is this section in five lines: what the corpus reports, what it stops short
of, what limits each conclusion, and the experiment that would settle it. The rest of the section gives
the evidence behind those five conclusions. Table~\ref{tab:statements} lists the \Fnrows{} evidence
statements in three classes, so that breadth is not read as strength: \Fnmulti{} reported across
multiple studies, \Fnsingle{} single-study or single-group signals and \Fncond{} conditional or
descriptive inventories. The inventories are rows defined by a selection: F7 collects the dictionaries
that attempted validation, and F9b, F10 and F11 list heterogeneous records. Every statement describes
this corpus, and none is a claim about the size or composition of the literature.

A few conventions apply to every number. A statement is a proposition that at least one record asserts,
and statements are split where records assert strictly different claims: F1b, F1b$'$ and F1b$''$
separate held-out recovery, gradient association and labelled correspondence, and F8a, F8a$'$ and
F8a$''$ separate an instrumented failure, a disagreement and an implausibility. In all, of the \Nclaims{}
records \Fnorow{} support no row and \Fmultirow{} are counted in more than one, \Fmultithree{} of them
in three. \Fflagrec{} of them, from \Fdefstud{} studies, carry a recorded defect of one of five
result-invalidating kinds (resubstitution, leakage across split units, circular estimation, selection on
the reported metric, or single-case fitting~\cite{kapoor2023leakage}) and are marked with $^{\dagger}$
wherever they appear rather than removed. Counts weight neither sample size nor effect size, and the same
datasets recur (CheXpert, MIMIC-CXR and TCGA run through F1a, F2, F2b and F3a alone); the splits,
external evaluation and uncertainty that the corpus reports are in \S\ref{sec:eval-compare}, and a
dictionary's reproducibility in \Nsaestab{} of \Nsae{}. Where records disagree, the disagreement is
stated rather than averaged.

How well a statement's evidence is safeguarded is a separate question from how much of it there is, and
no grade combines the two. Fig.~\ref{fig:findings} shows, for each statement, the share of its records
that report each of five groups of controls; the supplement lists the \Vnclasses{} underlying classes
with what each licenses and how many records report it. A control validates the measurement, not the proposition: two
controlled records are not a replication of a finding, and a share is comparable within a class, never
across classes. The map is mostly empty. A permutation null appears in \Vperm{} of \Vpermn{} records
where a probe could carry one and a randomly initialised encoder in \Vlrn{} of \Vlrnn{}; the control
task of Hewitt and Liang~\cite{hewitt2019designing} appears in none. A matched control is reported in
\Vspc{} of \Vspcn{} targeted interventions, and a comparison model in \Vcomp{} of \Vcompn{} geometry
records.

\textbf{What these counts are.} A count here is a count of records or of studies in this corpus, and a
$k/n$ is the share of the records whose design can carry a control that report one. Neither is an effect
size, a probability or an evidence score, and no grade combines them, so breadth and safeguards are read
on two separate axes. On the first, a statement's \emph{corpus-breadth tier} is broad for at least five
distinct first authors, at least two modalities and at least three peer-reviewed records, moderate for at
least two distinct first authors together with either at least three of them or at least two modalities,
and narrow otherwise; \Fbroad{} statements are broad, and the supplement gives the distinct first authors
and the model lineages behind each. On the second, \Fzeroclass{} statements report no control in any class
their designs could carry, and \Fhalfclass{} report each control in fewer than half of the records that
could carry it. The two axes do not track each other: \Fbroadzero{} statement is broad and reports no
control at all. Records are not independent either: \Fmultirow{} of the \Nclaims{} are counted in more than
one statement, \Fmultithree{} of them in three, and the same datasets and groups recur across statements
(\S\ref{sec:eval-compare}).

\textbf{How far the statements depend on the weaker records.} Five restrictions were applied to every
statement, each reported as a change in how much of the corpus stands behind it. On peer-reviewed records
alone, \Fsensrefn{} statements change: \Fsensrefemptyn{} lose every record (\Fsensrefempty{}) and
\Fsensreffalln{} fall to a narrower tier. On the \Nminvalid{} studies that state a patient-, slide- or
subject-level split, evaluate on data unused for fitting and report seeds, intervals or tests,
\Fsensminemptyn{} statements lose every record, \Fsensminfalln{} fall a tier and \Fsensminkeep{} is
unchanged. With the records flagged above for one of the five result-invalidating kinds set aside,
\Fscrdrop{} statements lose records and \Fscrtier{} change tier. Collapsing the records to one per
last-author group, a coarse proxy for dependence between studies, changes \Fsenslastn{} statement
(\Fsenslast{}), and collapsing to one per model lineage --- CLIP-style dual encoders, DINO-style
self-supervised encoders, multimodal LLMs, task-specific classifiers --- changes \Fsenslinn{}
(\Fsenslin{}). On the stratum that the recorded queries return on a re-run, \Frepdrop{} statements lose a
record: \Frepdetail{}; \Freptier{} of them change tier. One coding judgement recurs through all five, whether
a study with several applicable records counts as controlled when one of them reports a control or only
when all do; it moves the controlled count from \Fattribany{} to \Fattribcons{} of \Fattribstud{}
study--statement pairs.

\begin{table*}[!tb]
\centering\small\setlength\tabcolsep{3pt}\renewcommand{\arraystretch}{1.02}
\caption{The studies behind each of the \Fnrows{} evidence statements, in three classes, with the number of studies (Stud.) and the modalities they span (Mod.). Fig.~\ref{fig:findings} states each proposition in full and gives the controls its records report.}
\label{tab:statements}
\begin{tabular}{@{}lL{4.6cm}ccL{7.6cm}@{}}
\toprule
Row & Statement & Stud. & Mod. & Studies \\
\midrule
\multicolumn{5}{@{}l}{\emph{Reported across multiple studies}} \\
F1a & Findings and anatomy, linear readouts & 18 & 7 & {\scriptsize\cite{beyeler2024comparing,chevli2026frozenct,demir2025downstream,farquhar2026pretraining,germani2025bias,goetze2025prediction,jung2026foundation,kasireddy2025explainable,li2025featurequality,liang2026cheap,lin2026transparent,marza2025thunder,muthyala2026frozen,pedersen2026look,sadman2025interpreting,steiner2026when,svatko2026deep,zhou2025generalist}} \\
F1a$'$ & The same, non-linear readouts & 15 & 5 & {\scriptsize\cite{aerts2025foundation,akanova2026interpreting,cui2026translating,denner2024leveraging,gustafsson2024evaluating,howard2024generative,hu2026probing,jamzad2026p3ca,li2025embeddings,neidlinger2025benchmark,nooralahzadeh2026sparse,peng2025benchmarking,ramchandani2026benchmarking,tamura2025patch,vadori2025mind}} \\
F1b & Held-out recovery, text-anchored scorers & 6 & 4 & {\scriptsize\cite{bertolini2025strengths,kim2024transparent,luo2025xbench,mensah2026interpretable,nie2025explainable,patricio2023conceptbased}} \\
F1b$'$ & Concept-vector gradient association & 2 & 2 & {\scriptsize\cite{hieromnimon2025analysis,nicolson2024textcavs}} \\
F1b$^{\prime\prime}$ & Labelled representational correspondence & 5 & 3 & {\scriptsize\cite{arasteh2026candor,bakhta2026frozen,mishra2025comparing,vig2026do,zhao2026benchmarking}} \\
F1c & Concept layers on foundation features & 12 & 7 & {\scriptsize\cite{gao2024aligning,gao2025evaluating,han2026clear,patricio2025twostep,pellegrini2023xplainer,rojas2026echopcbm,sun2025label,tushar2026distilling,wang2025mvp,yan2023robust,yang2024textbook,zhang2026regionfm}} \\
F2 & Sex, age or race decodable & 9 & 5 & {\scriptsize\cite{cajas2026shortkit,encin2026comparative,kim2026encoding,li2026acquisition,lin2025beyond,sourget2025fairness,tariq2024self,weber2023posthoc,zheng2024demographic}} \\
F2b & Acquisition site and device decodable & 12 & 4 & {\scriptsize\cite{dejong2025current,hajishafiezahramini2026dataset,kmen2026robust,komen2024do,li2026acquisition,lin2025beyond,mayer2025lossy,murchan2024deep,rahbar2026frozen,saueressig2025histology,zhang2025hospital,zhang2026auditing}} \\
F3a & Centre and scanner, pathology & 17 & 1 & {\scriptsize\cite{carloni2025pathology,dejong2025current,esbri2025benchmarking,gustafsson2024evaluating,jung2026foundation,kmen2026robust,komen2024do,lin2025beyond,mayer2025lossy,mishra2025comparing,murchan2024deep,saueressig2025histology,shafique2026investigating,thiringer2026scanner,wolflein2023benchmarking,zhang2025hospital,zhang2026auditing}} \\
F3b & Harmonisation raises invariance & 5 & 2 & {\scriptsize\cite{carloni2025pathology,donle2026featmap,kmen2026robust,murchan2024deep,zhang2025hospital}} \\
F3c & Site and artefact, brain MRI & 3 & 1 & {\scriptsize\cite{mielcarz2026do,navarrogonzalez2026crossscanner,rahbar2026frozen}} \\
F5b & Shared structure; grouping by objective & 4 & 3 & {\scriptsize\cite{aerts2025foundation,arasteh2026candor,folco2025semantic,tayebiarasteh2026self}} \\
F8a & Attribution fails a faithfulness test & 6 & 4 & {\scriptsize\cite{akanova2026interpreting,bors2026benchmarking,idaji2026beyond,mcinerney2022that,sadanandan2026attention,xiong2026rethinking}} \\
F8a$'$ & Attribution methods disagree & 3 & 3 & {\scriptsize\cite{bors2026benchmarking,idaji2026beyond,injarabian2025interpreting}} \\
F8a$^{\prime\prime}$ & Maps implausible on inspection & 2 & 2 & {\scriptsize\cite{hashmi2024towards,injarabian2025interpreting}} \\
F9a & Targeted interventions change outputs & 8 & 4 & {\scriptsize\cite{bouzid2025insights,chen2026vihd,farina2026translating,nooralahzadeh2026sparse,nooralahzadeh2026universal,sadanandan2026mechanistically,sadashivaiah2024explaining,zeng2026easylens}} \\
\multicolumn{5}{@{}l}{\emph{Single-study or single-group signals}} \\
F4 & Sex redundant for a refitted probe & 1 & 1 & {\scriptsize\cite{kim2026encoding}} \\
F5 & Anisotropy and whitening & 1 & 1 & {\scriptsize\cite{garciahidalgo2026embedding}} \\
F6 & Modality gap; separation task-dependent & 1 & 3 & {\scriptsize\cite{restrepo2026cone}} \\
F8b & Labels less decodable after the connector & 1 & 1 & {\scriptsize\cite{zhu2026lost}} \\
\multicolumn{5}{@{}l}{\emph{Conditional or descriptive inventories}} \\
F7 & Validation of dictionary features & 12 & 7 & {\scriptsize\cite{bisson2026unsupervised,dasdelen2025cytosae,kim2026dissecting,le2024learning,nakka2025mammosae,nerrise2026geosae,redlich2026explainable,renzulli2025medsae,sun2025prototypebased,wesp2026sparse,yan2026a,zhao2026compositional}} \\
F9b & Inventory: editing and unlearning & 14 & 6 & {\scriptsize\cite{arora2026semantic,balogh2026regularization,carloni2025pathology,donle2026featmap,folco2025semantic,garciahidalgo2026embedding,kim2026encoding,kmen2026robust,li2026acquisition,murchan2024deep,restrepo2026cone,sadashivaiah2024explaining,weber2023posthoc,zhang2025hospital}} \\
F10 & Inventory: medical vs general decodability & 20 & 7 & {\scriptsize\cite{bakhta2026frozen,denner2024leveraging,farquhar2026pretraining,germani2025bias,han2026clear,kim2024transparent,li2025featurequality,luo2025xbench,marza2025thunder,mensah2026interpretable,nie2025explainable,patricio2023conceptbased,patricio2025twostep,peng2025benchmarking,svatko2026deep,vadori2025mind,vig2026do,yan2023robust,zhao2026benchmarking,zhou2025generalist}} \\
F11 & Inventory: properties shared or differing & 19 & 7 & {\scriptsize\cite{arasteh2026candor,bertolini2025strengths,cajas2026shortkit,farquhar2026pretraining,folco2025semantic,garciahidalgo2026embedding,germani2025bias,jin2024fairmedfm,kim2026encoding,li2026acquisition,mishra2025comparing,muthyala2026frozen,ouaari2026assessing,restrepo2026cone,sadanandan2026attention,tayebiarasteh2026crossmodal,tayebiarasteh2026self,xiong2026rethinking,zoellin2024block}} \\
\bottomrule
\end{tabular}
\end{table*}


The five subsections below give the evidence behind the five conclusions of
Table~\ref{tab:synthesis}, one each, in its order. Each names the statements of
Table~\ref{tab:statements} it rests on rather than only their row codes.

\subsection{Availability against use}

Five statements report that a concept is decodable: clinical findings and anatomy under linear and under
non-linear readouts, held-out recovery by text-anchored scorers, concept layers on foundation features,
and acquisition site and device. They rest on 18, 15, 6, 12 and 11 distinct first authors and on four or
more modalities each. The statements of most practical consequence rest on far less. That a decodable
attribute is linearly redundant for a refitted readout rests on one record; that attributions fail an
instrumented faithfulness test, on six studies, with method disagreement and qualitative implausibility
carried as separate, weaker rows; and that a targeted intervention exceeds a matched control, on three
records, only one of them against random directions. The interventional record is too heterogeneous and too weakly
controlled to support any claim about magnitude: the eight records to which a matched-direction control
applies use incomparable end-points, none reports a dose--response curve, and two vary the magnitude
while tuning it~\cite{chen2026vihd,nooralahzadeh2026universal}. The effects can be listed but not
compared, from under 0.02 AUROC for a post-hoc gap shift~\cite{restrepo2026cone} and under 0.001 for
residualising a demographic attribute~\cite{kim2026encoding}$^{\dagger}$ to the collapse of a finding's
own AUROC when the roughly ten channels that decode it are zeroed~\cite{nooralahzadeh2026sparse}.
Post-hoc operations on frozen embeddings --- harmonisation, whitening, gap shifting, erasure and
editing --- are the only interventions with more than a handful of records.

\subsection{Acquisition and site structure}

Centre, scanner and device are recoverable from pathology and brain-MRI embeddings in every record of
the three statements that test them --- site and device decodability, centre and scanner in pathology,
and site and artefact in brain MRI --- except one, by 11, 16 and 3 distinct first authors, and where both
are measured, centre rivals or exceeds biological structure. The benefit of removing it depends on the
operation. Among the harmonisation records, affine harmonisation raises cross-scanner retrieval and
centre removal raises a robustness index without improving downstream generalisation, and no record
reports the downstream end-point after harmonisation on a held-out site.

\subsection{Meaning of discovered features}

Naming and validation are separate steps, and an automatically
named feature can still be validated: of the \Nsaeval{} dictionaries with a reported semantic check,
\Nsaemanual{} are named manually, \Nsaelabmatch{} by label matching, \Nsaellm{} by a language model and
\Nsaetext{} by text matching. The validated dictionaries share one property, their locus: all read
encoder-side representations, and none of the three trained on decoder residual streams is checked
against clinicians or structured labels. The two sides also report different kinds of feature, which
suggests, as a hypothesis confounded by model, layer, dictionary design and naming, that
residual-stream dictionaries of report generators surface report-level rather than image-level
concepts. \Nsaestab{} of the \Nsae{} report any reproducibility check, and none establishes the
cross-seed identity of its features.

\subsection{Medical against general pre-training}

\Nmvgrec{} core records from \Nmvgstud{} studies compare a
medical with a general foundation model on the same data; seventeen further records compare only with
an ImageNet- or task-supervised classifier and are excluded. Two inventory rows hold them, one for the
decoding records and one for the records in which a property is shared or differs; three records are
marked not compared, because the
general model is measured only on zero-shot accuracy rather than on the coded property, and three with
a determinate outcome fall outside both rows. Where the medical model is better it is better by degree, under
linear probing, frozen-feature adaptation, expert concept scoring and pathology
segmentation~\cite{li2025featurequality,germani2025bias,zhou2025generalist,kim2024transparent,song2024morphological,gildenblat2024segmentation,cui2026translating};
against this, general or specialist encoders match or beat it across more than 200 radiology
evaluations, in radiological retrieval, among ten CT models, in complementary rather than better
pathology features, and after post-hoc
alignment~\cite{baharoon2023dinov2,denner2024leveraging,chevli2026frozenct,neidlinger2025benchmark,folco2025semantic,wesp2026sparse}.
The properties that matter most for interpretability are shared by both kinds of model, from low
effective dimensionality and the same degradation under whitening to shortcut reliance on the same
radiographs, the loss of small-lesion signal at global pooling, subgroup disparities, the modality gap
and nearest neighbours that carry the opposite
label~\cite{garciahidalgo2026embedding,li2025featurequality,sourget2025fairness,muthyala2026frozen,jin2024fairmedfm,restrepo2026cone,arasteh2026candor}. All of these are convenience comparisons that vary architecture,
objective, scale and data at once; the one study that stratifies the objective in a controlled matrix
does not vary medical against general pre-training~\cite{tayebiarasteh2026self}. The effect of medical
pre-training itself is therefore unidentified, and no study tests whether an interpretability method
transfers between the two kinds of model (\S\ref{sec:future}).

\subsection{Where to measure}

Varying the readout within one model, as three studies do, suggests \textemdash{}
as a hypothesis rather than a statement \textemdash{} that readout granularity can dominate apparent decodability for small or
spatially localised targets: pooled probes sit at chance where patch-local probes on the same forward
pass reach 0.899 AUC, although only under an oracle lesion box~\cite{muthyala2026frozen}; shallow blocks
carry spatial information better than deep ones~\cite{vadori2025mind}; and acquisition site occupies a
large principal component~\cite{rahbar2026frozen}. No design separates model, depth and pooling, so
what follows for practice is only that an availability claim should name its sublocus as well as its
model. Beyond the encoder, class labels become less decodable after the connector in one uncontrolled
successive-locus probe, and no other core record measures a medical connector or the text encoder
as a locus of its own.
